\documentclass{article}
\usepackage{graphicx} % Extended graphics inclusions
\usepackage{float}
\usepackage{url} % For \url{}
\usepackage{../config/atxy} % For front cover
\usepackage{amsfonts} % Needed for some fonts
\usepackage[usenames]{color} % Needed for colored R input/output
\usepackage{pdfcolmk} % Correct some problems with the color stack


\title{Test suite: run the don't run}
\author{Lobry, J.R.}

\usepackage{/Library/Frameworks/R.framework/Resources/share/texmf/Sweave}
\begin{document}
%
% To change the R input/output style:
%
\definecolor{Soutput}{rgb}{0,0,0.56}
\definecolor{Sinput}{rgb}{0.56,0,0}
\DefineVerbatimEnvironment{Sinput}{Verbatim}
{formatcom={\color{Sinput}},fontsize=\footnotesize, baselinestretch=0.75}
\DefineVerbatimEnvironment{Soutput}{Verbatim}
{formatcom={\color{Soutput}},fontsize=\footnotesize, baselinestretch=0.75}
%
% Rlogo
%
\newcommand{\Rlogo}{\protect\includegraphics[height=1.8ex,keepaspectratio]{../figs/Rlogo.pdf}}
%
% Shortcut for seqinR:
%
\newcommand{\seqinr}{\texttt{seqin\bf{R}}}
\newcommand{\Seqinr}{\texttt{Seqin\bf{R}}}
\fvset{fontsize= \scriptsize}
%
% R output options and libraries to be loaded.
%
%
%  Sweave Options
%
% Put all figures in the fig folder and start the name with current file name.
% Do not produce EPS files
%


\maketitle
\tableofcontents
% BEGIN - DO NOT REMOVE THIS LINE

\section{Introduction}

Many seqinR functions use socket connections to retrieve information from
the internet. As a consequence, most of examples should be protected by
a \verb!\dontrun{}! to pass the \texttt{R CMD CHECK}. In this section
we want to run automatically all these examples to check that everything
is OK.


\section{Stop list}

This is the list of function that don't run for now and need to be fixed.

\begin{Schunk}
\begin{Sinput}
 stoplist <- c("GetFromSequence", "reverse.align")
\end{Sinput}
\end{Schunk}

% GetFromSequence: see later
% reverse.align: need clustalw on line, see later

\section{Don't run generator}

This code chunk generates the \texttt{dontrun.rnw} file that is included there after.
This file should be pre-existent, and two \texttt{Sweave()} passes are necessary.

\begin{Schunk}
\begin{Sinput}
 outfile <- file(paste(pwd, "dontrun.rnw", sep = "/"), open = "w")
 fex <- dir()
 for (f in fex) {
     fctname <- substr(x = f, start = 1, stop = nchar(f) - 
         2)
     if (fctname %in% stoplist) 
         next
     lines <- readLines(f)
     dontrun <- lines[which(substring(lines, 1, 3) == "##D")]
     if (length(dontrun) == 0) 
         next
     dontrun <- sapply(dontrun, function(x) substr(x, 5, nchar(x)))
     writeLines(paste("\\subsection{\\texttt{", fctname, "()}}", 
         sep = ""), outfile)
     writeLines(paste("<<", fctname, ",fig=F,keep.source=T>>=", 
         sep = ""), outfile)
     writeLines(dontrun, outfile)
     writeLines("@", outfile)
 }
 close(outfile)
\end{Sinput}
\end{Schunk}

\subsection{\texttt{SeqAcnucWeb()}}
\begin{Schunk}
\begin{Sinput}
 # Need internet connection
   choosebank("emblTP")
   query("mylist", "sp=felis catus et t=cds et o=mitochondrion")
   stopifnot(is.SeqAcnucWeb(mylist$req[[1]]))
   closebank()
\end{Sinput}
\end{Schunk}
\subsection{\texttt{alllistranks()}}
\begin{Schunk}
\begin{Sinput}
 # Need internet connection
  choosebank("emblTP")
  query("tmp1", "sp=Borrelia burgdorferi", virtual = TRUE)
  query("tmp2", "sp=Borrelia burgdorferi", virtual = TRUE)
  query("tmp3", "sp=Borrelia burgdorferi", virtual = TRUE)
  (result <- alllistranks())
\end{Sinput}
\begin{Soutput}
$count
[1] 3

$ranks
[1] 2 3 4
\end{Soutput}
\begin{Sinput}
  stopifnot(result$count == 3)   # Three ACNUC lists
  stopifnot(result$ranks == 2:4) # Starting at rank 2
  #
  # Summay of current lists defined on the ACNUC server:
  #
  sapply(result$ranks, getliststate)
\end{Sinput}
\begin{Soutput}
      [,1]   [,2]   [,3]  
type  "SQ"   "SQ"   "SQ"  
name  "TMP1" "TMP2" "TMP3"
count 1682   1682   1682  
locus TRUE   TRUE   TRUE  
\end{Soutput}
\begin{Sinput}
  closebank()
\end{Sinput}
\end{Schunk}
\subsection{\texttt{autosocket()}}
\begin{Schunk}
\begin{Sinput}
  #Need internet connection
   choosebank("emblTP")
   autosocket()
\end{Sinput}
\begin{Soutput}
                description                       class 
"->pbil.univ-lyon1.fr:5558"                    "socket" 
                       mode                        text 
                       "a+"                      "text" 
                     opened                    can read 
                   "opened"                       "yes" 
                  can write 
                      "yes" 
\end{Soutput}
\begin{Sinput}
   closebank()
\end{Sinput}
\end{Schunk}
\subsection{\texttt{countsubseqs()}}
\begin{Schunk}
\begin{Sinput}
  # Need internet connection
   choosebank("emblTP")
   query("mylist", "N=@", virtual = TRUE) # select all (seqs + subseqs)
   mylist$nelem   # 14138094 seqs + subseqs
\end{Sinput}
\begin{Soutput}
[1] 14138094
\end{Soutput}
\begin{Sinput}
   stopifnot(mylist$nelem == 14138094)
   css(glr("mylist")) # 1604500 subsequences only
\end{Sinput}
\begin{Soutput}
[1] 1604500
\end{Soutput}
\begin{Sinput}
   stopifnot(css(glr("mylist")) == 1604500)
   closebank()
\end{Sinput}
\end{Schunk}
\subsection{\texttt{crelistfromclientdata()}}
\begin{Schunk}
\begin{Sinput}
  # Need internet connection
  choosebank("emblTP")
  #
  # Example with a file that contains sequence names:
  #
  fileSQ <- system.file("sequences/bb.mne", package = "seqinr")
  crelistfromclientdata("listSQ", file = fileSQ, type = "SQ")
  sapply(listSQ$req, getName)
\end{Sinput}
\begin{Soutput}
 [1] "A04009.OSPA"   "A04009.OSPB"   "A22442"        "A24006"       
 [5] "A24008"        "A24010"        "A24012"        "A24014"       
 [9] "A24016"        "A33362"        "A67759.PE1"    "AB011063"     
[13] "AB011064"      "AB011065"      "AB011066"      "AB011067"     
[17] "AB035616"      "AB035617"      "AB035618"      "AB041949.VLSE"
\end{Soutput}
\begin{Sinput}
  #
  # Example with a file that contains sequence accession numbers:
  #
  fileAC <- system.file("sequences/bb.acc", package = "seqinr")
  crelistfromclientdata("listAC", file = fileAC, type = "AC")
  sapply(listAC$req, getName) 
\end{Sinput}
\begin{Soutput}
 [1] "AY382159" "AY382160" "AY491412" "AY498719" "AY498720" "AY498721"
 [7] "AY498722" "AY498723" "AY498724" "AY498725" "AY498726" "AY498727"
[13] "AY498728" "AY498729" "AY499181" "AY500379" "AY500380" "AY500381"
[19] "AY500382" "AY500383"
\end{Soutput}
\begin{Sinput}
  #
  # Example with a file that contains species names:
  #
  fileSP <- system.file("sequences/bb.sp", package = "seqinr")
  crelistfromclientdata("listSP", file = fileSP, type = "SP")
  sapply(listSP$req, getName) 
\end{Sinput}
\begin{Soutput}
 [1] "BORRELIA ANSERINA"    "BORRELIA CORIACEAE"   "BORRELIA PARKERI"    
 [4] "BORRELIA TURICATAE"   "BORRELIA HERMSII"     "BORRELIA CROCIDURAE" 
 [7] "BORRELIA LONESTARI"   "BORRELIA HISPANICA"   "BORRELIA BARBOURI"   
[10] "BORRELIA THEILERI"    "BORRELIA DUTTONII"    "BORRELIA MIYAMOTOI"  
[13] "BORRELIA PERSICA"     "BORRELIA RECURRENTIS" "BORRELIA BURGDORFERI"
[16] "BORRELIA AFZELII"     "BORRELIA GARINII"     "BORRELIA ANDERSONII" 
[19] "BORRELIA VALAISIANA"  "BORRELIA JAPONICA"   
\end{Soutput}
\begin{Sinput}
  #
  # Example with a file that contains keywords:
  #
  fileKW <- system.file("sequences/bb.kwd", package = "seqinr")
  crelistfromclientdata("listKW", file = fileKW, type = "KW")
  sapply(listKW$req, getName)
\end{Sinput}
\begin{Soutput}
 [1] "PLASMID"         "CIRCULAR"        "PARTIAL"         "5'-PARTIAL"     
 [5] "3'-PARTIAL"      "MOTA GENE"       "MOTB GENE"       "DIVISION PRO"   
 [9] "GYRB GENE"       "JOINING REGION"  "FTSA GENE"       "RPOB GENE"      
[13] "RPOC GENE"       "FLA GENE"        "DNAJ GENE"       "TUF GENE"       
[17] "PGK GENE"        "RUVA GENE"       "RUVB GENE"       "PROMOTER REGION"
\end{Soutput}
\begin{Sinput}
  #
  # Summary of ACNUC lists:
  #
  sapply(alr()$rank, getliststate)
\end{Sinput}
\begin{Soutput}
      [,1]          [,2]     [,3]          [,4]     [,5]          [,6]    
type  "SQ"          "SQ"     "SQ"          "SQ"     "SP"          "SP"    
name  "LOCALFILE_1" "LISTSQ" "LOCALFILE_2" "LISTAC" "LOCALFILE_3" "LISTSP"
count 20            20       20            20       20            20      
locus FALSE         FALSE    TRUE          TRUE     TRUE          TRUE    
      [,7]          [,8]    
type  "KW"          "KW"    
name  "LOCALFILE_4" "LISTKW"
count 20            20      
locus TRUE          TRUE    
\end{Soutput}
\begin{Sinput}
  closebank() 
\end{Sinput}
\end{Schunk}
\subsection{\texttt{draw.rearranged.oriloc()}}
\begin{Schunk}
\begin{Sinput}
 r.ori <- rearranged.oriloc(seq.fasta = system.file("sequences/ct.fasta",package = "seqinr"),
     g2.coord = system.file("sequences/ct.coord",package = "seqinr"))
\end{Sinput}
\end{Schunk}
\subsection{\texttt{extract.breakpoints()}}
\begin{Schunk}
\begin{Sinput}
 r.ori <- rearranged.oriloc(seq.fasta = system.file("sequences/ct.fasta",package = "seqinr"),
     g2.coord = system.file("sequences/ct.coord",package = "seqinr"))
\end{Sinput}
\end{Schunk}
\subsection{\texttt{extractseqs()}}
\begin{Schunk}
\begin{Sinput}
 # Need internet connection
  choosebank("swissprotTP")
  query("mylist", "k=globin", virtual = TRUE)
  mylist.fasta <- exseq("mylist", verbose = TRUE)
\end{Sinput}
\begin{Soutput}
I'm checking the arguments...
... and everything is OK up to now.
Format is  fasta 
Operation is  simple 
The rank of the list  mylist is  2 .
request :  extractseqs&lrank=2&format=fasta&operation=simple&zlib=T 
Running getzlibsock... 
'con' is a connection...
Socket number is 5....
Trying to get answer from socket...
n=1000, nn=1000,nnn=1000
extractseqs successfully ended ...
Number of lines     : 102
Number of sequences : 28
extractseqs OK, program carry on...
Ok, everything is fine!
Closing socket close_sock_gz_r  status = 0
\end{Soutput}
\begin{Sinput}
  # 103 lines of FASTA 
  stopifnot(length(mylist.fasta) == 103)
  closebank()
\end{Sinput}
\end{Schunk}
\subsection{\texttt{getType()}}
\begin{Schunk}
\begin{Sinput}
 # Need internet connection
   choosebank("emblTP")
   getType()
\end{Sinput}
\begin{Soutput}
        sname                                  libel
2661      CDS              .PE protein coding region
2662       ID                            Locus entry
2663 MISC_RNA .RN other structural RNA coding region
2664     RRNA          .RR Ribosomal RNA coding gene
2665    SCRNA              .SC small cytoplasmic RNA
2666    SNRNA                  .SN small nuclear RNA
2667     TRNA           .TR Transfer RNA coding gene
\end{Soutput}
\end{Schunk}
\subsection{\texttt{getlistrank()}}
\begin{Schunk}
\begin{Sinput}
  # Need internet connection
  choosebank("emblTP")
  query("MyListName", "sp=Borrelia burgdorferi", virtual = TRUE)
  (result <- getlistrank("MyListName"))
\end{Sinput}
\begin{Soutput}
[1] 2
\end{Soutput}
\begin{Sinput}
  stopifnot(result == 2)
  closebank()
\end{Sinput}
\end{Schunk}
\subsection{\texttt{getliststate()}}
\begin{Schunk}
\begin{Sinput}
   ### Need internet connection
   choosebank("emblTP")
   query("mylist", "sp=felis catus et t=cds", virtual=TRUE)
   getliststate(glr("mylist")) # SQ, MYLIST, 603, FALSE
\end{Sinput}
\begin{Soutput}
$type
[1] "SQ"

$name
[1] "MYLIST"

$count
[1] 603

$locus
[1] FALSE
\end{Soutput}
\begin{Sinput}
   gln(glr("mylist")) # MYLIST (upper case letters on server)
\end{Sinput}
\begin{Soutput}
[1] "MYLIST"
\end{Soutput}
\begin{Sinput}
   closebank()
\end{Sinput}
\end{Schunk}
\subsection{\texttt{gfrag()}}
\begin{Schunk}
\begin{Sinput}
 # Need internet connection
   choosebank("emblTP")
   gfrag("LMFLCHR36", start = 1, length = 3529852) -> myseq
   stopifnot(nchar(myseq) == 3529852)
   closebank()   
\end{Sinput}
\end{Schunk}
\subsection{\texttt{ghelp()}}
\begin{Schunk}
\begin{Sinput}
   ### Need internet connection
   choosebank("emblTP")
   ghelp()
\end{Sinput}
\begin{Soutput}
      ---- General Information on ACNUC nucleic acid data base ----             
HELP:                                                                           
   A detailed explanation of purpose and usage of each command is obtained      
by typing the command name and requesting help when the dialog suggests it.     
SEQUENCES AND SUBSEQUENCES:                                                     
   In addition to sequences as published in research articles, ACNUC contains   
subsequences which are sequence segments with specific coding function (e.g.    
protein, tRNA, rRNA genes...). Sequence type distinguishes parent from sub-     
sequences: parent sequences have ID type, subsequences have a type that      
indicates their function (CDS, TRNA, RRNA,...). Most subsequence names derive   
from the parent sequence's name by addition of suffixes .PEn, .TRn, .RRn, .SNn  
.RNn for CDS, TRNA, RRNA, snRNA or misc_RNA-typed subsequences, respectively.   
When the gene name is known, it is used as a suffix in the corresponding        
subsequence name.
SEQUENCE LISTS:                                                                 
   This program deals with sequence lists which group sequences selected from   
the data base using one or more selection criteria (see SELECT help). Many      
sequence lists can be handled simultaneously by the program and previous lists  
can be used to define new ones.                                                 
   Typical use of program is:                                                   
 - SPECIES command to know which species names are to be used in selection.     
 - KEYWORDS command to know which keywords are to be used in selection.         
 - SELECT command to select sequences from data base combining various          
criteria. This command produces the list of sequences that fit the criteria.    
 - SHORT command to obtain a brief description of selected sequences            
or - INFO command to get more detailed information.                             
 - EXTRACT command to copy selected sequences to a user file.                   
LIST NAMES:                                                                     
   Lists are created by commands SELECT or FIND. They are given automatically a 
name (LIST1, LIST2,...) by the program, unless the user enters his own list     
name by appending /l=my_list_name to the command name at the "Command?" prompt. 
Most commands operate either on a sequence list or on an individual sequence.   
Reply to question "List, sequence, or accession #? [default=...]" with          
<RETURN> to access the default (list of) sequence(s) or with any list name,     
sequence name, or accession number.                                             
FILE OUTPUT:                                                                    
If /lpt is appended to command name at the "Command?" prompt, the output of     
commands SPECIES, KEYWORDS, INFO, SHORT, NAMES, CODES, BASES goes to a file     
named `query.out'.                                                              
CODED NAMES:                                                                    
   Coded names are to be used when specifying species, keywords, journals,      
sequence types, organelles, molecules. Specific commands    
(SPECIES, KEYWORDS, CODES) allow you to find these names easily.                
REFERENCES:                                                                     
   To find a sequence from a bibliographical reference use the selection        
criterion "R=reference-code" of SELECT command. Build the reference code as     
follows (journal names are given by CODES command):                             
  journal_name/volume/first_page         for journal articles                   
  book/year/name_of_1st_author           for books                              
  thesis/year/name_of_1st_author         for thesis                             
  patent/patent_number                   for patented sequences                 
  unpubl/year/name_of_1st_author         for unpublished sequences              
Example: nar/8/2173 stands for Nucleic Acids Research 8:2173-2192 (1980).       
\end{Soutput}
\begin{Sinput}
   ghelp("SELECT")
\end{Sinput}
\begin{Soutput}
     In addition to functions described in the help for the simple usage of     
command SELECT, other selection criteria and operations between lists exist.    
Specifically, it is also possible to build lists of species and lists of        
keywords for further retrieval capabilities.                                    
                                                                                
Criteria      Resulting selection                                               
FK=file name  List of keywords taken from a file (which may have been created   
              by a SAVE command).                                               
FS=file name  List of species taken from a file (which may have been created    
              by a SAVE command).                                               
                                                                                
Operation     Result                                                            
ME list       Replaces subsequences in list by sequences from which they        
              are extracted (equivalent to option 4 of command MODIFY).         
FI list       Sequences in list plus all of their subsequences (equivalent to   
              option 5 of command MODIFY).                                      
PS list       Produces the list of species names attached to sequences in list. 
PK list       Produces the list of keyword names attached to sequences in list. 
UN list       If applied to a species list, produces the list of sequences from 
              species in the list; if applied to a keyword list, produces the   
              list of sequences attached to keywords in list.                   
SD spec-list  Applied to a list of species, produces the list of all descendants
              from them in the species tree. The list itself can easily be      
              created by command FIND.                                          
KD keyw-list  Applied to a list of keywords, produces the list of all           
              descendants from them in the keywords tree. The list itself can   
              easily be created by command FIND.                                
                                                                                
Operators PS, PK, and UN allow to solve the problem "find all genes             
simultaneously sequenced in a given series of species".                         
First, build the lists of sequences from each of these species. Next project    
each of these lists to attached keywords by applying operator PK. Then compute  
the list of keywords in common by combining the keyword lists with operator ET. 
Then, remove from this list of common keywords, those which are                 
uncharacteristic (e.g. partial) by employing command MODIFY. Finally, produce   
the lists of sequences attached to common keywords from each species by         
applying operator UN combined with initial species-based sequence lists.        
Species and keyword lists can be listed with command NAMES and saved with SAVE. 
\end{Soutput}
\begin{Sinput}
   # To get info about current database:
   ghelp("CONT")
\end{Sinput}
\begin{Soutput}
             ****     ACNUC Data Base Content      ****                         
              EMBL Library Release 78 WITHOUT ESTs  (March 2004)
27,571,397,913 bases; 12,533,594 sequences; 1,604,500 subseqs; 339,186 refers.
Software by M. Gouy & M. Jacobzone, Laboratoire de biometrie, Universite Lyon I 
\end{Soutput}
\end{Schunk}
\subsection{\texttt{isenum()}}
\begin{Schunk}
\begin{Sinput}
   ### Need internet connection
   choosebank("emblTP")
   isenum("LMFLCHR36")
\end{Sinput}
\begin{Soutput}
$number
[1] 13682678

$length
[1] 3529852

$frame
[1] 0

$gencode
[1] 0

$ncbigc
[1] 1

$otheraccessmatches
[1] FALSE
\end{Soutput}
\begin{Sinput}
   isn("LMFLCHR36")
\end{Sinput}
\begin{Soutput}
[1] 13682678
\end{Soutput}
\begin{Sinput}
   stopifnot(isn("LMFLCHR36") == 13682678)
   # Example with CDS:
   isenum("AB004237")
\end{Sinput}
\begin{Soutput}
$number
[1] 66351

$length
[1] 1140

$frame
[1] 0

$gencode
[1] 2

$ncbigc
[1] 2

$otheraccessmatches
[1] FALSE
\end{Soutput}
\end{Schunk}
\subsection{\texttt{knowndbs()}}
\begin{Schunk}
\begin{Sinput}
   ### Need internet connection
   choosebank("emblTP")
   kdb()
\end{Sinput}
\begin{Soutput}
           bank status
1       genbank     on
2          embl     on
3       emblwgs     on
4     swissprot     on
5       ensembl     on
6        refseq     on
7     hobacnucl     on
8     hobacprot     on
9     hovernucl     on
10    hoverprot     on
11    hogennucl     on
12    hogenprot     on
13   hogen4nucl     on
14   hogen4prot     on
15 homolensprot     on
16 homolensnucl     on
17      greview     on
18    HAMAPnucl     on
19    HAMAPprot     on
20    hoppsigen     on
21    nurebnucl     on
22    nurebprot     on
23   taxobacgen     on
                                                                            info
1                   GenBank Rel. 161 (15 August 2007) Last Updated: Oct 21, 2007
2            EMBL Library Release 92 (September 2007) Last Updated: Oct 18, 2007
3               EMBL Whole Genome Shotgun sequences Release 92  (September 2007)
4        UniProt Rel. 12 (SWISS-PROT 54 + TrEMBL 37): Last Updated: Oct  8, 2007
5                                                   Ensembl databases release 41
6                        RefSeq 15.0 (1 January 2006) Last Updated: Jan 23, 2006
7                        HOBACGEN - genomic data - Release 10 (February 12 2002)
8                        HOBACGEN - protein data - Release 10 (February 12 2002)
9  HOVERGEN - genomic data - Release 48 (May 24 2007) Last Updated: May 24, 2007
10 HOVERGEN - protein data - Release 48 (May 24 2007) Last Updated: May 24, 2007
11  HOGENOM - genomic data - Release 03 (Oct 14 2005) Last Updated: Nov  7, 2005
12  HOGENOM - protein data - Release 03 (Oct 14 2005) Last Updated: Mar 10, 2006
13 HOGENOM - genomic data - Release 04 (Sept 18,2007) Last Updated: Oct  2, 2007
14 HOGENOM - protein data - Release 04 (Sept 18,2007) Last Updated: Sep 18, 2007
15         HOMOLENS 3 - Homologous genes from Ensembl Last Updated: Jan 19, 2007
16        HOMOLENS 3 Homologous genes from Ensembl 41 Last Updated: Jan 19, 2007
17               EBI Genome Reviews. Acnuc Release 7. Last Updated: Feb 26, 2007
18              HAMAP - Acnuc Release - Nucleotides - Last Updated: Jun  5, 2007
19                   HAMAP - Acnuc Release -Proteins- Last Updated: Jun  5, 2007
20                                                                     Hoppsigen
21                   Nurebase 4.0 (26 September 2003) Last Updated: NOV 27, 2003
22                   Nurebase 4.0 (26 September 2003) Last Updated: NOV 27, 2003
23                                            TaxoBacGen Rel. 7 (September 2005)
\end{Soutput}
\begin{Sinput}
   closebank()
\end{Sinput}
\end{Schunk}
\subsection{\texttt{modifylist()}}
\begin{Schunk}
\begin{Sinput}
   ### Need internet connection
   choosebank("emblTP")
   query("mylist", "sp=felis catus et t=cds", virtual=TRUE)
   mylist$nelem # 603 sequences
\end{Sinput}
\begin{Soutput}
[1] 603
\end{Soutput}
\begin{Sinput}
   # select sequences with at least 1000 bp
   modifylist("mylist", operation = ">1000", virtual = TRUE)
   mylist.mod$nelem # now, only 132 sequences
\end{Sinput}
\begin{Soutput}
[1] 132
\end{Soutput}
\begin{Sinput}
   # scan for "felis" in annotations:
   modifylist("mylist.mod", "MyFinalList", op = "felis", type = "scan", virt =T)
   MyFinalList$nelem # now, only 33 sequences
\end{Sinput}
\begin{Soutput}
[1] 33
\end{Soutput}
\begin{Sinput}
   # Summary of current ACNUC lists:
   sapply(alr()$rank, getliststate)
\end{Sinput}
\begin{Soutput}
      [,1]     [,2]         [,3]         
type  "SQ"     "SQ"         "SQ"         
name  "MYLIST" "MYLIST.MOD" "MYFINALLIST"
count 603      132          33           
locus FALSE    FALSE        FALSE        
\end{Soutput}
\end{Schunk}
\subsection{\texttt{oriloc()}}
\begin{Schunk}
\begin{Sinput}
 #
 # A little bit too long for routine checks because oriloc() is already
 # called in draw.oriloc.Rd documentation file. Try example(draw.oriloc)
 # instead, or copy/paste the following code:
 #
 out <- oriloc()
 plot(out$st, out$sk, type = "l", xlab = "Map position in Kb",
     ylab = "Cumulated composite skew", 
     main = expression(italic(Chlamydia~~trachomatis)~~complete~~genome))
\end{Sinput}
\end{Schunk}
\subsection{\texttt{plot.SeqAcnucWeb()}}
\begin{Schunk}
\begin{Sinput}
   ### Need internet connection
   choosebank("hovernucl")
   query("list", "AC=AB000425")  
   plot(list$req[[1]])
\end{Sinput}
\end{Schunk}
\subsection{\texttt{prettyseq()}}
\begin{Schunk}
\begin{Sinput}
   ### Need internet connection
   choosebank("emblTP")
   prettyseq(111)
\end{Sinput}
\begin{Soutput}
Name: A00165           Length:108
Genetic code used: NUG=AUN=M when initiation codon

                  10         20         30         40         50         60 
           Q  Y  C   G  N  L  S   T  C  M   L  G  T   Y  T  Q  D   F  N  K
          cagtactgcg gtaatctgag tacttgcatg ctgggcacat acacgcagga cttcaacaag 
          >A00165

                  70         80         90        100        110 
           F  H  T   F  P  Q  T   A  I  G   V  G  A   P  G  *
          tttcacacgt tcccccaaac tgcaattggg gttggagcac ctggttga   
                                                       A00165<
\end{Soutput}
\end{Schunk}
\subsection{\texttt{print.SeqAcnucWeb()}}
\begin{Schunk}
\begin{Sinput}
   ### Need internet connection
   choosebank("emblTP")
   query("mylist", "sp=felis catus")     
   mylist$req[[1]]
\end{Sinput}
\begin{Soutput}
    name   length    frame   ncbicg 
"A06937"     "34"      "0"      "1" 
\end{Soutput}
\end{Schunk}
\subsection{\texttt{query()}}
\begin{Schunk}
\begin{Sinput}
  # Need internet connection
  choosebank("genbank")
  query("bb", "sp=Borrelia burgdorferi")
  # To get the names of the 4 first sequences:
  sapply(bb$req[1:4], getName)
\end{Sinput}
\begin{Soutput}
[1] "A04009" "A22442" "A24006" "A24008"
\end{Soutput}
\begin{Sinput}
  # To get the 4 first sequences:
  sapply(bb$req[1:4], getSequence, as.string = TRUE) 
\end{Sinput}
\begin{Soutput}
[1] "aagcttaattagaaccaaacttaattaaaaccaaacttaattgaagttattatcattttattttttttcaattttctatttgttatttgttaatcttataatataattatacttgtattaagttatattaatataaaaggagaatatattatgaaaaaatatttattgggaataggtctaatattagccttaatagcatgtaagcaaaatgttagcagccttgacgagaaaaacagcgtttcagtagatttgcctggtgaaatgaaagttcttgtaagcaaagaaaaaaacaaagacggcaagtacgatctaattgcaacagtagacaagcttgagcttaaaggaacttctgataaaaacaatggatctggagtacttgaaggcgtaaaagctgacaaaagtaaagtaaaattaacaatttctgacgatctaggtcaaaccacacttgaagttttcaaagaagatggcaaaacactagtatcaaaaaaagtaacttccaaagacaagtcatcaacagaagaaaaattcaatgaaaaaggtgaagtatctgaaaaaataataacaagagcagacggaaccagacttgaatacacaggaattaaaagcgatggatctggaaaagctaaagaggttttaaaaggctatgttcttgaaggaactctaactgctgaaaaaacaacattggtggttaaagaaggaactgttactttaagcaaaaatatttcaaaatctggggaagtttcagttgaacttaatgacactgacagtagtgctgctactaaaaaaactgcagcttggaattcaggcacttcaactttaacaattactgtaaacagtaaaaaaactaaagaccttgtgtttacaaaagaaaacacaattacagtacaacaatacgactcaaatggcaccaaattagaggggtcagcagttgaaattacaaaacttgatgaaattaaaaacgctttaaaataaggagaatttatgagattattaataggatttgctttagcgttagctttaataggatgtgcacaaaaaggtgctgagtcaattggttctcaaaaagaaaatgatctaaaccttgaagactctagtaaaaaatcacatcaaaacgctaaacaagaccttcctgcggtgacagaagactcagtgtctttgtttaatggtaataaaatttttgtaagcaaagaaaaaaatagctccggcaaatatgatttaagagcaacaattgatcaggttgaacttaaaggaacttccgataaaaacaatggttctggaacccttgaaggttcaaagcctgacaagagtaaagtaaaattaacagtttctgctgatttaaacacagtaaccttagaagcatttgatgccagcaaccaaaaaatttcaagtaaagttactaaaaaacaggggtcaataacagaggaaactctcaaagctaataaattagactcaaagaaattaacaagatcaaacggaactacacttgaatactcacaaataacagatgctgacaatgctacaaaagcagtagaaactctaaaaaatagcattaagcttgaaggaagtcttgtagtcggaaaaacaacagtggaaattaaagaaggtactgttactctaaaaagagaaattgaaaaagatggaaaagtaaaagtctttttgaatgacactgcaggttctaacaaaaaaacaggtaaatgggaagacagtactagcactttaacaattagtgctgacagcaaaaaaactaaagatttggtgttcttaacagatggtacaattacagtacaacaatacaacacagctggaaccagcctagaaggatcagcaagtgaaattaaaaatctttcagagcttaaaaacgctttaaaataatatataagtaaaccccctacaaggcatcagctagtgtaggaag"
[2] "atgaaaaaatatttattgggaataggtctaatattagccttaatagcatgtaagcaaaatgttagcagccttgacgagaaaaacagcgtttcagtagatttgcctggtgaaatgaacgttcttgtaagcaaagaaaaaaacaaagacggcaagtacgatctaattgcaacagtagacaagcttgagcttaaaggaacttctgataaaaacaatggatctggagtacttgaaggcgtaaaagctgacaaaagtaaagtaaaattaacaatttctgacgatctaggtcaaaccacacttgaagttttcaaagaagatggcaaaacactagtatcaaaaaaagtaacttccaaagacaagtcatcaacagaagaaaaattcaatgaaaaaggtgaagtatctgaaaaaataataacaagagcagacggaaccagacttgaatacacagaaattaaaagcgatggatctggaaaagctaaagaggttttaaaaagctatgttcttgaaggaactttaactgctgaaaaaacaacattggtggttaaagaaggaactgttactttaagcaaaaatatttcaaaatctggggaagtttcagttgaacttaatgacactgacagtagtgctgctactaaaaaaactgcagcttggaattcaggcacttcaactttaacaattactgtaaacagtaaaaaaactaaagaccttgtgtttacaaaagaaaacacaattacagtacaacaatacgactcaaatggcaccaaattagaggggtcagcagttgaaattacaaaacttgatgaaattaaaaacgctttaaaataa"                                                                                                                                                                                                                                                                                                                                                                                                                                                                                                                                                                                                                                                                                                                                                                                                                                                                                                                                                                                                                                                                                                                                     
[3] "atgaaaaaatatttattgggaataggtctaatattagccttaatagcatgtaagcaaaatgttagcagccttgatgaaaaaaatagcgtttcagtagatttacctggtggaatgaaagttcttgtaagtaaagaaaaagacaaagatggtaaatacagtctagaggcaacagtagacaagcttgagcttaaaggaacttctgataaaaacaacggttctggaacacttgaaggtgaaaaaactgacaaaagtaaagtaaaattaacaattgctgaggatctaagtaaaaccacatttgaaattttcaaagaagatggcaaaacattagtatcaaaaaaagtaacccttaaagacaagtcatcaacagaagaaaaattcaacgaaaagggtgaaatatctgaaaaaacaatagtaagagcaaatggaaccagacttgaatacacagacataaaaagcgatggatccggaaaagctaaagaagttttaaaagactttactcttgaaggaactctagctgctgacggcaaaacaacattgaaagttacagaaggcactgttgttttaagcaagaacattttaaaatccggagaaataacagttgcacttgatgactctgacactactcaggctactaaaaaaactggaaaatgggattcaaagacttccactttaacaattagtgtgaatagccaaaaaaccaaaaaccttgtattcacaaaagaagacacaataacagtacaaaaatacgactcagcaggcaccaatctagaaggcaaagcagtcgaaattacaacacttaaagaacttaaagacgctttaaaataa"                                                                                                                                                                                                                                                                                                                                                                                                                                                                                                                                                                                                                                                                                                                                                                                                                                                                                                                                                                                                                                                                                                                                  
[4] "atgaaaaaatatttattgggaataggtctaatattagccttaatagcatgtaagcaaaatgttagcagccttgacgagaaaaacagcgtttcagtagatgtacctggtggaatgaaagttcttgtaagcaaagaaaaaaacaaagacggcaagtacgatctaatggcaacagtggacaacgttgatcttaaaggaacttctgacaaaaacaatggatctggaatacttgaaggcgtaaaagctgataaaagtaaagtaaaattaacagttgctgacgatctaagcaaaaccacacttgaagttttaaaagaagatggtacagtagtgtcaagaaaagtaacttccaaagacaagtcaacaacagaagcaaaattcaacgaaaaaggtgaattgtctgaaaaaacaatgacaagagcaaacggaactactcttgaatactcacaaatgacaaatgaagacaatgctgcaaaagcagtagaaactcttaaaaacggcattaagtttgaaggaaatctcgcaagtggaaaaacagcagtggaaattaaagaaggcactgttactctaaaaagagaaattgataaaaatggaaaagtaaccgtctctttaaatgacactgcatctggttctaaaaaaacagcttcctggcaagaaagtactagcaccttaacaattagtgcaaacagcaaaaaaactaaagatctagtgttcctaacaaacggtacaattacagtacaaaattatgactcagctggcactaaacttgaaggatcagcagctgaaattaaaaaactcgatgaacttaaaaacgctttaagataa"                                                                                                                                                                                                                                                                                                                                                                                                                                                                                                                                                                                                                                                                                                                                                                                                                                                                                                                                                                                                                                                                                                                                        
\end{Soutput}
\end{Schunk}
\subsection{\texttt{readfirstrec()}}
\begin{Schunk}
\begin{Sinput}
 # Need internet connection
   choosebank("genbank")
   allowedtype <- readfirstrec()
   sapply(allowedtype, function(x) readfirstrec(type = x))
\end{Sinput}
\begin{Soutput}
      AUT       BIB       ACC       SMJ       SUB       LOC       KEY 
   164084    481229  77706747      4889  82184766  78700468   7277619 
     SPEC      SHRT       LNG       EXT       TXT 
   495418 706070711  22581444   7242727    400402 
\end{Soutput}
\end{Schunk}
\subsection{\texttt{rearranged.oriloc()}}
\begin{Schunk}
\begin{Sinput}
 r.ori <- rearranged.oriloc(seq.fasta = system.file("sequences/ct.fasta",package = "seqinr"),
     g2.coord = system.file("sequences/ct.coord",package = "seqinr"))
\end{Sinput}
\end{Schunk}
\subsection{\texttt{residuecount()}}
\begin{Schunk}
\begin{Sinput}
   ### Need internet connection
   choosebank("emblTP")
   query("mylist", "t=CDS", virtual = TRUE)
   stopifnot(residuecount(glr("mylist")) == 1611439240)
   stopifnot(is.na(residuecount(glr("unknowlist")))) # A warning is issued
\end{Sinput}
\end{Schunk}
\subsection{\texttt{savelist()}}
\begin{Schunk}
\begin{Sinput}
   ### Need internet connection
   choosebank("emblTP")
   query("mylist", "sp=felis catus et t=cds", virtual=TRUE)
   savelist(glr("mylist"))
\end{Sinput}
\begin{Soutput}
603 sequence mnemonics written into file: MYLIST.mne
\end{Soutput}
\begin{Sinput}
   # 603 sequence mnemonics written into file: MYLIST.mne 
   savelist(glr("mylist"), type = "A")
\end{Sinput}
\begin{Soutput}
603 sequence accession numbers written into file: MYLIST.acc
\end{Soutput}
\end{Schunk}
\subsection{\texttt{setlistname()}}
\begin{Schunk}
\begin{Sinput}
   ### Need internet connection
   choosebank("emblTP")
   query("mylist", "sp=felis catus et t=CDS", virtual = TRUE)
   # Change list name on server:
   setlistname(lrank = glr("mylist"), name = "feliscatus") # 0, OK.
\end{Sinput}
\begin{Soutput}
[1] 0
\end{Soutput}
\begin{Sinput}
   glr("mylist") # 0, list doesn't exist no more.
\end{Sinput}
\begin{Soutput}
[1] 0
\end{Soutput}
\begin{Sinput}
   glr("feliscatus") # 2, this list exists.
\end{Sinput}
\begin{Soutput}
[1] 2
\end{Soutput}
\end{Schunk}
\subsection{\texttt{translate()}}
\begin{Schunk}
\begin{Sinput}
 ## Need internet connection.
 ## Translation of the following EMBL entry:
 ##
 ## FT   CDS             join(complement(153944..154157),complement(153727..153866),
 ## FT                   complement(152185..153037),138523..138735,138795..138955)
 ## FT                   /codon_start=1
 ## FT                   /db_xref="FLYBASE:FBgn0002781"
 ## FT                   /db_xref="GOA:Q86B86"
 ## FT                   /db_xref="TrEMBL:Q86B86"
 ## FT                   /note="mod(mdg4) gene product from transcript CG32491-RZ;
 ## FT                   trans splicing"
 ## FT                   /gene="mod(mdg4)"
 ## FT                   /product="CG32491-PZ"
 ## FT                   /locus_tag="CG32491"
 ## FT                   /protein_id="AAO41581.1"
 ## FT                   /translation="MADDEQFSLCWNNFNTNLSAGFHESLCRGDLVDVSLAAEGQIVKA
 ## FT                   HRLVLSVCSPFFRKMFTQMPSNTHAIVFLNNVSHSALKDLIQFMYCGEVNVKQDALPAF
 ## FT                   ISTAESLQIKGLTDNDPAPQPPQESSPPPAAPHVQQQQIPAQRVQRQQPRASARYKIET
 ## FT                   VDDGLGDEKQSTTQIVIQTTAAPQATIVQQQQPQQAAQQIQSQQLQTGTTTTATLVSTN
 ## FT                   KRSAQRSSLTPASSSAGVKRSKTSTSANVMDPLDSTTETGATTTAQLVPQQITVQTSVV
 ## FT                   SAAEAKLHQQSPQQVRQEEAEYIDLPMELPTKSEPDYSEDHGDAAGDAEGTYVEDDTYG
 ## FT                   DMRYDDSYFTENEDAGNQTAANTSGGGVTATTSKAVVKQQSQNYSESSFVDTSGDQGNT
 ## FT                   EAQVTQHVRNCGPQMFLISRKGGTLLTINNFVYRSNLKFFGKSNNILYWECVQNRSVKC
 ## FT                   RSRLKTIGDDLYVTNDVHNHMGDNKRIEAAKAAGMLIHKKLSSLTAADKIQGSWKMDTE
 ## FT                   GNPDHLPKM"
 choosebank("emblTP")
 query("trans", "N=AE003734.PE35")
 getTrans(trans$req[[1]])
\end{Sinput}
\begin{Soutput}
  [1] "M" "A" "D" "D" "E" "Q" "F" "S" "L" "C" "W" "N" "N" "F" "N" "T" "N" "L"
 [19] "S" "A" "G" "F" "H" "E" "S" "L" "C" "R" "G" "D" "L" "V" "D" "V" "S" "L"
 [37] "A" "A" "E" "G" "Q" "I" "V" "K" "A" "H" "R" "L" "V" "L" "S" "V" "C" "S"
 [55] "P" "F" "F" "R" "K" "M" "F" "T" "Q" "M" "P" "S" "N" "T" "H" "A" "I" "V"
 [73] "F" "L" "N" "N" "V" "S" "H" "S" "A" "L" "K" "D" "L" "I" "Q" "F" "M" "Y"
 [91] "C" "G" "E" "V" "N" "V" "K" "Q" "D" "A" "L" "P" "A" "F" "I" "S" "T" "A"
[109] "E" "S" "L" "Q" "I" "K" "G" "L" "T" "D" "N" "D" "P" "A" "P" "Q" "P" "P"
[127] "Q" "E" "S" "S" "P" "P" "P" "A" "A" "P" "H" "V" "Q" "Q" "Q" "Q" "I" "P"
[145] "A" "Q" "R" "V" "Q" "R" "Q" "Q" "P" "R" "A" "S" "A" "R" "Y" "K" "I" "E"
[163] "T" "V" "D" "D" "G" "L" "G" "D" "E" "K" "Q" "S" "T" "T" "Q" "I" "V" "I"
[181] "Q" "T" "T" "A" "A" "P" "Q" "A" "T" "I" "V" "Q" "Q" "Q" "Q" "P" "Q" "Q"
[199] "A" "A" "Q" "Q" "I" "Q" "S" "Q" "Q" "L" "Q" "T" "G" "T" "T" "T" "T" "A"
[217] "T" "L" "V" "S" "T" "N" "K" "R" "S" "A" "Q" "R" "S" "S" "L" "T" "P" "A"
[235] "S" "S" "S" "A" "G" "V" "K" "R" "S" "K" "T" "S" "T" "S" "A" "N" "V" "M"
[253] "D" "P" "L" "D" "S" "T" "T" "E" "T" "G" "A" "T" "T" "T" "A" "Q" "L" "V"
[271] "P" "Q" "Q" "I" "T" "V" "Q" "T" "S" "V" "V" "S" "A" "A" "E" "A" "K" "L"
[289] "H" "Q" "Q" "S" "P" "Q" "Q" "V" "R" "Q" "E" "E" "A" "E" "Y" "I" "D" "L"
[307] "P" "M" "E" "L" "P" "T" "K" "S" "E" "P" "D" "Y" "S" "E" "D" "H" "G" "D"
[325] "A" "A" "G" "D" "A" "E" "G" "T" "Y" "V" "E" "D" "D" "T" "Y" "G" "D" "M"
[343] "R" "Y" "D" "D" "S" "Y" "F" "T" "E" "N" "E" "D" "A" "G" "N" "Q" "T" "A"
[361] "A" "N" "T" "S" "G" "G" "G" "V" "T" "A" "T" "T" "S" "K" "A" "V" "V" "K"
[379] "Q" "Q" "S" "Q" "N" "Y" "S" "E" "S" "S" "F" "V" "D" "T" "S" "G" "D" "Q"
[397] "G" "N" "T" "E" "A" "Q" "V" "T" "Q" "H" "V" "R" "N" "C" "G" "P" "Q" "M"
[415] "F" "L" "I" "S" "R" "K" "G" "G" "T" "L" "L" "T" "I" "N" "N" "F" "V" "Y"
[433] "R" "S" "N" "L" "K" "F" "F" "G" "K" "S" "N" "N" "I" "L" "Y" "W" "E" "C"
[451] "V" "Q" "N" "R" "S" "V" "K" "C" "R" "S" "R" "L" "K" "T" "I" "G" "D" "D"
[469] "L" "Y" "V" "T" "N" "D" "V" "H" "N" "H" "M" "G" "D" "N" "K" "R" "I" "E"
[487] "A" "A" "K" "A" "A" "G" "M" "L" "I" "H" "K" "K" "L" "S" "S" "L" "T" "A"
[505] "A" "D" "K" "I" "Q" "G" "S" "W" "K" "M" "D" "T" "E" "G" "N" "P" "D" "H"
[523] "L" "P" "K" "M" "*"
\end{Soutput}
\end{Schunk}

\begin{Schunk}
\begin{Sinput}
 setwd(pwd)
\end{Sinput}
\end{Schunk}

\section{Session Informations}

This part was compiled under the following \Rlogo{}~environment:

\begin{itemize}
  \item R version 2.6.0 (2007-10-03), \verb|i386-apple-darwin8.10.1|
  \item Locale: \verb|C|
  \item Base packages: base, datasets, grDevices, graphics, methods,
    stats, utils
  \item Other packages: MASS~7.2-36, ade4~1.4-4, ape~2.0-1,
    gee~4.13-13, lattice~0.16-5, nlme~3.1-85, seqinr~1.1-3,
    xtable~1.5-1
  \item Loaded via a namespace (and not attached): grid~2.6.0,
    rcompgen~0.1-15
\end{itemize}
There were two compilation steps:

\begin{itemize}
  \item \Rlogo{} compilation time was: Sun Oct 21 20:28:26 2007
  \item \LaTeX{} compilation time was: \today
\end{itemize}

% END - DO NOT REMOVE THIS LINE

%%%%%%%%%%%%  BIBLIOGRAPHY %%%%%%%%%%%%%%%%%
\clearpage
\addcontentsline{toc}{section}{References}
\bibliographystyle{plain}
\bibliography{../config/book}
\end{document}
